\documentclass[11pt]{article}
\newcommand{\videoid}{8tnjx0tEn4g}
% Preámbulo compartido de las notas del curso de topología algebraica.
% Cada nota empieza con \documentclass[11pt]{article}, define \videoid{...}
% (el identificador del video de YouTube) y luego hace % Preámbulo compartido de las notas del curso de topología algebraica.
% Cada nota empieza con \documentclass[11pt]{article}, define \videoid{...}
% (el identificador del video de YouTube) y luego hace % Preámbulo compartido de las notas del curso de topología algebraica.
% Cada nota empieza con \documentclass[11pt]{article}, define \videoid{...}
% (el identificador del video de YouTube) y luego hace \input{preambulo}.
% Compilar con LuaLaTeX:  lualatex nombre-de-la-nota.tex
\usepackage[spanish,provide=*]{babel}
\usepackage[T1]{fontenc}
\usepackage{lmodern}
\usepackage[a4paper,margin=2.6cm]{geometry}
\usepackage{amsmath,amssymb,amsthm}
\usepackage{xcolor}
\usepackage{tikz}
\usetikzlibrary{arrows.meta,babel}
\usepackage[colorlinks=true,urlcolor=blue!50!black,linkcolor=blue!50!black]{hyperref}

\hyphenpenalty=10000 \exhyphenpenalty=10000 \emergencystretch=3em

% Colores de las figuras (los mismos que usa la página)
\definecolor{extremo}{RGB}{47,107,79}
\definecolor{azul}{RGB}{44,74,154}

% Definiciones numeradas por unidad
\theoremstyle{definition}
\newtheorem{defn}{Definición}
% Numeración por unidad: \unidad{u}{n} (después de \input{preambulo}) hace que las
% definiciones de la nota se llamen u.n, u.(n+1), ...; n es el primer número
% libre de esa unidad (ver la tabla de CONVENCIONES.md).
\newcommand{\unidad}[2]{\renewcommand{\thedefn}{#1.\arabic{defn}}\setcounter{defn}{\numexpr#2-1\relax}}

% Observaciones y erratas del video. Usar SOLO si hay algo que señalar:
%   \begin{erratas}
%     \errata{203}{3:23}{Se dice ..., pero debe decir ...}
%   \end{erratas}
% El primer argumento son los segundos (enlaza al minuto exacto del video),
% el segundo es la marca visible, el tercero es el texto.
\newenvironment{erratas}{\section*{Observaciones y erratas del video}\begin{itemize}}{\end{itemize}}
\newcommand{\errata}[3]{%
  \edef\errurl{https://www.youtube.com/watch?v=\videoid\string&t=#1}%
  \item[\textbf{\expandafter\href\expandafter{\errurl}{#2}}] #3}
.
% Compilar con LuaLaTeX:  lualatex nombre-de-la-nota.tex
\usepackage[spanish,provide=*]{babel}
\usepackage[T1]{fontenc}
\usepackage{lmodern}
\usepackage[a4paper,margin=2.6cm]{geometry}
\usepackage{amsmath,amssymb,amsthm}
\usepackage{xcolor}
\usepackage{tikz}
\usetikzlibrary{arrows.meta,babel}
\usepackage[colorlinks=true,urlcolor=blue!50!black,linkcolor=blue!50!black]{hyperref}

\hyphenpenalty=10000 \exhyphenpenalty=10000 \emergencystretch=3em

% Colores de las figuras (los mismos que usa la página)
\definecolor{extremo}{RGB}{47,107,79}
\definecolor{azul}{RGB}{44,74,154}

% Definiciones numeradas por unidad
\theoremstyle{definition}
\newtheorem{defn}{Definición}
% Numeración por unidad: \unidad{u}{n} (después de % Preámbulo compartido de las notas del curso de topología algebraica.
% Cada nota empieza con \documentclass[11pt]{article}, define \videoid{...}
% (el identificador del video de YouTube) y luego hace \input{preambulo}.
% Compilar con LuaLaTeX:  lualatex nombre-de-la-nota.tex
\usepackage[spanish,provide=*]{babel}
\usepackage[T1]{fontenc}
\usepackage{lmodern}
\usepackage[a4paper,margin=2.6cm]{geometry}
\usepackage{amsmath,amssymb,amsthm}
\usepackage{xcolor}
\usepackage{tikz}
\usetikzlibrary{arrows.meta,babel}
\usepackage[colorlinks=true,urlcolor=blue!50!black,linkcolor=blue!50!black]{hyperref}

\hyphenpenalty=10000 \exhyphenpenalty=10000 \emergencystretch=3em

% Colores de las figuras (los mismos que usa la página)
\definecolor{extremo}{RGB}{47,107,79}
\definecolor{azul}{RGB}{44,74,154}

% Definiciones numeradas por unidad
\theoremstyle{definition}
\newtheorem{defn}{Definición}
% Numeración por unidad: \unidad{u}{n} (después de \input{preambulo}) hace que las
% definiciones de la nota se llamen u.n, u.(n+1), ...; n es el primer número
% libre de esa unidad (ver la tabla de CONVENCIONES.md).
\newcommand{\unidad}[2]{\renewcommand{\thedefn}{#1.\arabic{defn}}\setcounter{defn}{\numexpr#2-1\relax}}

% Observaciones y erratas del video. Usar SOLO si hay algo que señalar:
%   \begin{erratas}
%     \errata{203}{3:23}{Se dice ..., pero debe decir ...}
%   \end{erratas}
% El primer argumento son los segundos (enlaza al minuto exacto del video),
% el segundo es la marca visible, el tercero es el texto.
\newenvironment{erratas}{\section*{Observaciones y erratas del video}\begin{itemize}}{\end{itemize}}
\newcommand{\errata}[3]{%
  \edef\errurl{https://www.youtube.com/watch?v=\videoid\string&t=#1}%
  \item[\textbf{\expandafter\href\expandafter{\errurl}{#2}}] #3}
) hace que las
% definiciones de la nota se llamen u.n, u.(n+1), ...; n es el primer número
% libre de esa unidad (ver la tabla de CONVENCIONES.md).
\newcommand{\unidad}[2]{\renewcommand{\thedefn}{#1.\arabic{defn}}\setcounter{defn}{\numexpr#2-1\relax}}

% Observaciones y erratas del video. Usar SOLO si hay algo que señalar:
%   \begin{erratas}
%     \errata{203}{3:23}{Se dice ..., pero debe decir ...}
%   \end{erratas}
% El primer argumento son los segundos (enlaza al minuto exacto del video),
% el segundo es la marca visible, el tercero es el texto.
\newenvironment{erratas}{\section*{Observaciones y erratas del video}\begin{itemize}}{\end{itemize}}
\newcommand{\errata}[3]{%
  \edef\errurl{https://www.youtube.com/watch?v=\videoid\string&t=#1}%
  \item[\textbf{\expandafter\href\expandafter{\errurl}{#2}}] #3}
.
% Compilar con LuaLaTeX:  lualatex nombre-de-la-nota.tex
\usepackage[spanish,provide=*]{babel}
\usepackage[T1]{fontenc}
\usepackage{lmodern}
\usepackage[a4paper,margin=2.6cm]{geometry}
\usepackage{amsmath,amssymb,amsthm}
\usepackage{xcolor}
\usepackage{tikz}
\usetikzlibrary{arrows.meta,babel}
\usepackage[colorlinks=true,urlcolor=blue!50!black,linkcolor=blue!50!black]{hyperref}

\hyphenpenalty=10000 \exhyphenpenalty=10000 \emergencystretch=3em

% Colores de las figuras (los mismos que usa la página)
\definecolor{extremo}{RGB}{47,107,79}
\definecolor{azul}{RGB}{44,74,154}

% Definiciones numeradas por unidad
\theoremstyle{definition}
\newtheorem{defn}{Definición}
% Numeración por unidad: \unidad{u}{n} (después de % Preámbulo compartido de las notas del curso de topología algebraica.
% Cada nota empieza con \documentclass[11pt]{article}, define \videoid{...}
% (el identificador del video de YouTube) y luego hace % Preámbulo compartido de las notas del curso de topología algebraica.
% Cada nota empieza con \documentclass[11pt]{article}, define \videoid{...}
% (el identificador del video de YouTube) y luego hace \input{preambulo}.
% Compilar con LuaLaTeX:  lualatex nombre-de-la-nota.tex
\usepackage[spanish,provide=*]{babel}
\usepackage[T1]{fontenc}
\usepackage{lmodern}
\usepackage[a4paper,margin=2.6cm]{geometry}
\usepackage{amsmath,amssymb,amsthm}
\usepackage{xcolor}
\usepackage{tikz}
\usetikzlibrary{arrows.meta,babel}
\usepackage[colorlinks=true,urlcolor=blue!50!black,linkcolor=blue!50!black]{hyperref}

\hyphenpenalty=10000 \exhyphenpenalty=10000 \emergencystretch=3em

% Colores de las figuras (los mismos que usa la página)
\definecolor{extremo}{RGB}{47,107,79}
\definecolor{azul}{RGB}{44,74,154}

% Definiciones numeradas por unidad
\theoremstyle{definition}
\newtheorem{defn}{Definición}
% Numeración por unidad: \unidad{u}{n} (después de \input{preambulo}) hace que las
% definiciones de la nota se llamen u.n, u.(n+1), ...; n es el primer número
% libre de esa unidad (ver la tabla de CONVENCIONES.md).
\newcommand{\unidad}[2]{\renewcommand{\thedefn}{#1.\arabic{defn}}\setcounter{defn}{\numexpr#2-1\relax}}

% Observaciones y erratas del video. Usar SOLO si hay algo que señalar:
%   \begin{erratas}
%     \errata{203}{3:23}{Se dice ..., pero debe decir ...}
%   \end{erratas}
% El primer argumento son los segundos (enlaza al minuto exacto del video),
% el segundo es la marca visible, el tercero es el texto.
\newenvironment{erratas}{\section*{Observaciones y erratas del video}\begin{itemize}}{\end{itemize}}
\newcommand{\errata}[3]{%
  \edef\errurl{https://www.youtube.com/watch?v=\videoid\string&t=#1}%
  \item[\textbf{\expandafter\href\expandafter{\errurl}{#2}}] #3}
.
% Compilar con LuaLaTeX:  lualatex nombre-de-la-nota.tex
\usepackage[spanish,provide=*]{babel}
\usepackage[T1]{fontenc}
\usepackage{lmodern}
\usepackage[a4paper,margin=2.6cm]{geometry}
\usepackage{amsmath,amssymb,amsthm}
\usepackage{xcolor}
\usepackage{tikz}
\usetikzlibrary{arrows.meta,babel}
\usepackage[colorlinks=true,urlcolor=blue!50!black,linkcolor=blue!50!black]{hyperref}

\hyphenpenalty=10000 \exhyphenpenalty=10000 \emergencystretch=3em

% Colores de las figuras (los mismos que usa la página)
\definecolor{extremo}{RGB}{47,107,79}
\definecolor{azul}{RGB}{44,74,154}

% Definiciones numeradas por unidad
\theoremstyle{definition}
\newtheorem{defn}{Definición}
% Numeración por unidad: \unidad{u}{n} (después de % Preámbulo compartido de las notas del curso de topología algebraica.
% Cada nota empieza con \documentclass[11pt]{article}, define \videoid{...}
% (el identificador del video de YouTube) y luego hace \input{preambulo}.
% Compilar con LuaLaTeX:  lualatex nombre-de-la-nota.tex
\usepackage[spanish,provide=*]{babel}
\usepackage[T1]{fontenc}
\usepackage{lmodern}
\usepackage[a4paper,margin=2.6cm]{geometry}
\usepackage{amsmath,amssymb,amsthm}
\usepackage{xcolor}
\usepackage{tikz}
\usetikzlibrary{arrows.meta,babel}
\usepackage[colorlinks=true,urlcolor=blue!50!black,linkcolor=blue!50!black]{hyperref}

\hyphenpenalty=10000 \exhyphenpenalty=10000 \emergencystretch=3em

% Colores de las figuras (los mismos que usa la página)
\definecolor{extremo}{RGB}{47,107,79}
\definecolor{azul}{RGB}{44,74,154}

% Definiciones numeradas por unidad
\theoremstyle{definition}
\newtheorem{defn}{Definición}
% Numeración por unidad: \unidad{u}{n} (después de \input{preambulo}) hace que las
% definiciones de la nota se llamen u.n, u.(n+1), ...; n es el primer número
% libre de esa unidad (ver la tabla de CONVENCIONES.md).
\newcommand{\unidad}[2]{\renewcommand{\thedefn}{#1.\arabic{defn}}\setcounter{defn}{\numexpr#2-1\relax}}

% Observaciones y erratas del video. Usar SOLO si hay algo que señalar:
%   \begin{erratas}
%     \errata{203}{3:23}{Se dice ..., pero debe decir ...}
%   \end{erratas}
% El primer argumento son los segundos (enlaza al minuto exacto del video),
% el segundo es la marca visible, el tercero es el texto.
\newenvironment{erratas}{\section*{Observaciones y erratas del video}\begin{itemize}}{\end{itemize}}
\newcommand{\errata}[3]{%
  \edef\errurl{https://www.youtube.com/watch?v=\videoid\string&t=#1}%
  \item[\textbf{\expandafter\href\expandafter{\errurl}{#2}}] #3}
) hace que las
% definiciones de la nota se llamen u.n, u.(n+1), ...; n es el primer número
% libre de esa unidad (ver la tabla de CONVENCIONES.md).
\newcommand{\unidad}[2]{\renewcommand{\thedefn}{#1.\arabic{defn}}\setcounter{defn}{\numexpr#2-1\relax}}

% Observaciones y erratas del video. Usar SOLO si hay algo que señalar:
%   \begin{erratas}
%     \errata{203}{3:23}{Se dice ..., pero debe decir ...}
%   \end{erratas}
% El primer argumento son los segundos (enlaza al minuto exacto del video),
% el segundo es la marca visible, el tercero es el texto.
\newenvironment{erratas}{\section*{Observaciones y erratas del video}\begin{itemize}}{\end{itemize}}
\newcommand{\errata}[3]{%
  \edef\errurl{https://www.youtube.com/watch?v=\videoid\string&t=#1}%
  \item[\textbf{\expandafter\href\expandafter{\errurl}{#2}}] #3}
) hace que las
% definiciones de la nota se llamen u.n, u.(n+1), ...; n es el primer número
% libre de esa unidad (ver la tabla de CONVENCIONES.md).
\newcommand{\unidad}[2]{\renewcommand{\thedefn}{#1.\arabic{defn}}\setcounter{defn}{\numexpr#2-1\relax}}

% Observaciones y erratas del video. Usar SOLO si hay algo que señalar:
%   \begin{erratas}
%     \errata{203}{3:23}{Se dice ..., pero debe decir ...}
%   \end{erratas}
% El primer argumento son los segundos (enlaza al minuto exacto del video),
% el segundo es la marca visible, el tercero es el texto.
\newenvironment{erratas}{\section*{Observaciones y erratas del video}\begin{itemize}}{\end{itemize}}
\newcommand{\errata}[3]{%
  \edef\errurl{https://www.youtube.com/watch?v=\videoid\string&t=#1}%
  \item[\textbf{\expandafter\href\expandafter{\errurl}{#2}}] #3}

\unidad{1}{12} % unidad 1; primer enunciado: 1.12

\title{Cambio de punto base}
\author{Luis Jorge Sánchez Saldaña\\[2pt]\small Facultad de Ciencias, UNAM}
\date{}

\begin{document}
\maketitle

\noindent\small Notas del video
\href{https://www.youtube.com/watch?v=8tnjx0tEn4g}{\guillemotleft{}Cambio de punto base para el grupo fundamental\guillemotright{}}
(curso de Introducción a la Topología Algebraica). Estudian qué pasa con el grupo fundamental al cambiar el punto base: puede cambiar si los puntos no se pueden unir por un camino, pero si se pueden unir los grupos son isomorfos.
\normalsize

\section*{La pregunta}

Hasta ahora se tiene un espacio topológico cualquiera $X$, fijo, y un punto arbitrario pero fijo $x_0\in X$. El grupo fundamental de $X$ basado en $x_0$ es el conjunto de lazos basados en $x_0$ módulo la homotopía relativa a los extremos,
\[
  \pi_1(X,x_0)=\Omega(X,x_0)\big/\!\simeq\ \mathrm{rel}\ \{0,1\}.
\]
Lo primero que se observa en esta definición es que el grupo fundamental no solo depende de $X$, sino también del punto base. Entonces surge la pregunta natural:
\[
  ?`\pi_1(X,x_0)=\pi_1(X,x_1)\ \text{para otro punto } x_1\in X?
\]
Es decir, si se cambia de punto base, ?`se obtiene otro grupo o el mismo? Esa es la pregunta que se contesta en este video.

\section*{Un ejemplo: el punto base sí importa}

Sea $X=S^1\sqcup\{y\}$, la unión disjunta de un círculo y un punto, y sean $x_0\in S^1$ y $x_1=y$ como puntos base. Entonces
\[
  \pi_1(X,x_0)=\pi_1(S^1,x_0)\neq 1,
  \qquad
  \pi_1(X,x_1)=\pi_1(\{y\},y)=1.
\]
El video no demuestra todavía que $\pi_1(S^1,x_0)$ es no trivial (lo anuncia para los próximos videos, y pide creerlo por ahora); en cambio, el grupo fundamental de un punto se verifica fácilmente: solo hay un lazo, el trivial. Por lo tanto estos dos grupos no son iguales (ni isomorfos), y el grupo fundamental sí depende del punto base.

La razón es que el grupo fundamental solo reconoce lo que pasa en la componente conexa, y en realidad arco-conexa (esto se verá con más detalle después), en la que vive el punto base. Aquí los dos puntos base están en componentes distintas, y cada grupo solo \guillemotleft{}ve\guillemotright{} la suya.

\medskip\noindent\textbf{Respuesta parcial.} Si los puntos base viven en componentes arco-conexas diferentes, es decir, si no se pueden conectar por un camino, entonces los grupos fundamentales podrían ser diferentes.

Queda la otra situación: ?`qué pasa si los dos puntos base sí se pueden conectar por un camino? Aquí está la respuesta.

\section*{El teorema}

\begin{teo}[Cambio de punto base]
Si $x_0,x_1\in X$ se pueden conectar por un camino, entonces
\[
  \pi_1(X,x_0)\cong\pi_1(X,x_1).
\]
\end{teo}

Antes de demostrarlo, conviene pensar geométricamente qué está pasando. Se dibujan $X$, los puntos $x_0$ y $x_1$ y un camino $h$ que los conecta. Se quiere definir un morfismo de grupos entre $\pi_1(X,x_0)$ y $\pi_1(X,x_1)$; ni siquiera es claro cómo definir una función entre ellos. Un elemento de $\pi_1(X,x_0)$ es la clase de un lazo $f$ basado en $x_0$, y se necesita un lazo basado en $x_1$. La idea es la siguiente: se comienza en $x_1$, se recorre $h$ hasta llegar a $x_0$, se recorre $f$ y luego se recorre $h$ al revés para regresar a $x_1$. Aquí $\bar h$ denota $h$ parametrizado en la dirección contraria: $\bar h(s)=h(1-s)$.

\begin{center}
\begin{tikzpicture}[x=0.0165cm,y=-0.0165cm,line cap=round]
  \draw[gray,rounded corners=12pt] (10,10) rectangle (470,210);
  \node at (448,34) {$X$};
  \draw[very thick] (120,140) .. controls (90,60) and (30,100) .. (45,150) .. controls (55,185) and (105,185) .. (120,140);
  \node at (22,118) {$f$};
  \draw[very thick,azul,-{Stealth[length=7pt]}] (390,60) .. controls (320,30) and (260,130) .. (120,140);
  \node[azul] at (292,66) {$h$};
  \fill[extremo] (120,140) circle (4);
  \fill[extremo] (390,60) circle (4);
  \node at (140,162) {$x_0$};
  \node at (404,44) {$x_1$};
\end{tikzpicture}
\end{center}

\noindent Figura: el camino $h$ va de $x_1$ a $x_0$ y $f$ es un lazo basado en $x_0$. El lazo $h\cdot f\cdot\bar h$ comienza en $x_1$, recorre $h$, luego $f$ y regresa a $x_1$ por $\bar h$.

\medskip\noindent\textbf{Observación sobre el sentido de $h$.} Para que la concatenación $h\cdot f\cdot\bar h$ tenga sentido, $h$ debe ser un camino que comienza en $x_1$ y termina en $x_0$: así $h$ termina donde comienza $f$, y $\bar h$ comienza donde termina $f$. En lo que sigue, $h$ es un camino de $x_1$ a $x_0$ (si $x_0$ y $x_1$ se conectan por un camino en un sentido, se conectan en el otro por el camino inverso).

\medskip\noindent\textbf{Comentario.} Como en el video, se escribe $h\cdot f\cdot\bar h$ sin paréntesis. Cualquier manera de poner los paréntesis da la misma clase de homotopía relativa a los extremos, porque la concatenación es asociativa salvo homotopía; la homotopía que se usó para la asociatividad en el video anterior (Teorema~1.11) funciona igual para caminos que no son lazos. Del mismo modo, $\bar h\cdot h\simeq c_{x_0}$ y $h\cdot\bar h\simeq c_{x_1}$ rel $\{0,1\}$ se obtienen con la homotopía de los inversos de ese video, que solo usa que $h$ es un camino.

\begin{defn}[cambio de punto base]
Sea $h$ un camino de $x_1$ a $x_0$. Se define
\[
  \beta_h\colon \pi_1(X,x_0)\longrightarrow\pi_1(X,x_1),
  \qquad
  [f]\longmapsto[h\cdot f\cdot\bar h].
\]
\end{defn}

\begin{proof}
Se prueba que $\beta_h$ es un isomorfismo. Para concluir hay que verificar tres cosas: (1) $\beta_h$ está bien definida; (2) $\beta_h$ es un homomorfismo de grupos; (3) $\beta_h$ es biyectiva.

\medskip\noindent\textit{(1) $\beta_h$ está bien definida.}
Hay que ver que el resultado no depende del representante elegido de la clase de $f$. Sea $f'$ otro lazo basado en $x_0$ con $f'\simeq f\ \mathrm{rel}\ \{0,1\}$. Como la concatenación respeta la homotopía relativa a los extremos (Proposición~1.8, vista en un video anterior), se obtiene
\[
  h\cdot f'\cdot\bar h\simeq h\cdot f\cdot\bar h\ \mathrm{rel}\ \{0,1\},
  \qquad\text{es decir,}\qquad
  [h\cdot f'\cdot\bar h]=[h\cdot f\cdot\bar h].
\]
En el pizarrón esto se escribe de manera más general: si $h'\simeq h$, $\bar h'\simeq\bar h$ y $f'\simeq f$ rel $\{0,1\}$, entonces $h'\cdot f'\cdot\bar h'\simeq h\cdot f\cdot\bar h$ rel $\{0,1\}$. Así, la clase no depende de los representantes elegidos ni de nada más, y $\beta_h$ está bien definida como función.

\medskip\noindent\textit{(2) $\beta_h$ es un homomorfismo.}
Sean $[f],[g]\in\pi_1(X,x_0)$. Por definición de $\beta_h$,
\[
  \beta_h([f][g])=\beta_h[f\cdot g]=[h\cdot f\cdot g\cdot\bar h].
\]
Ahora se hace el viejo truco de multiplicar por $\bar h\cdot h$ en medio:
\[
  [h\cdot f\cdot g\cdot\bar h]
  =[h\cdot f\cdot\bar h\cdot h\cdot g\cdot\bar h]
  =[h\cdot f\cdot\bar h]\,[h\cdot g\cdot\bar h]
  =\beta_h[f]\,\beta_h[g].
\]
?`Por qué es legal el truco? Porque es muy geométrico. El lazo $h\cdot f\cdot g\cdot\bar h$ recorre $h$, luego $f$, luego $g$ y finalmente $h$ al revés. El otro lazo, $h\cdot f\cdot\bar h\cdot h\cdot g\cdot\bar h$, recorre $h$, $f$, luego $h$ hacia atrás, luego $h$ hacia adelante, luego $g$ y de nuevo $h$ hacia atrás. Cuando se recorre $\bar h$ y luego $h$ se forma una especie de lacito en medio, que se puede ir recogiendo y contraer hasta el punto $x_0$. Es decir, $\bar h\cdot h\simeq c_{x_0}$ rel $\{0,1\}$, donde $c_{x_0}$ es el lazo constante en $x_0$, y entonces los dos lazos son homotópicos relativamente a los extremos; basta contraer el pedacito de en medio. Esto da la primera igualdad. La segunda es la definición del producto en $\pi_1(X,x_1)$.

\medskip\noindent\textit{(3) $\beta_h$ es biyectiva.}
Para esto se da una inversa. Se define el mismo morfismo pero ahora con el camino $\bar h$, que va en sentido contrario ($\bar h$ va de $x_0$ a $x_1$):
\[
  \beta_{\bar h}\colon\pi_1(X,x_1)\longrightarrow\pi_1(X,x_0),
  \qquad
  [f]\longmapsto[\bar h\cdot f\cdot h].
\]
Se afirma que $\beta_{\bar h}$ es inversa de $\beta_h$. Para $[f]\in\pi_1(X,x_1)$,
\[
  \beta_h\beta_{\bar h}[f]
  =\beta_h[\bar h\cdot f\cdot h]
  =[h\cdot\bar h\cdot f\cdot h\cdot\bar h]
  =[f],
\]
porque $h\cdot\bar h$ y $c_{x_1}$ son homotópicos relativamente a los extremos, igual que en el dibujo de hace rato (el video invita a convencerse geométricamente). Análogamente se demuestra que $\beta_{\bar h}\beta_h[f]=[f]$ para $[f]\in\pi_1(X,x_0)$. Por lo tanto $\beta_h$ y $\beta_{\bar h}$ son inversas una de la otra, y $\beta_h$ es biyectiva.

Con esto $\beta_h$ es un homomorfismo biyectivo, es decir, un isomorfismo entre el grupo fundamental de $X$ basado en $x_0$ y el basado en $x_1$.
\end{proof}

\begin{erratas}
  \errata{217}{3:37}{Se dice que $h$ es un camino que comienza en $x_0$ y termina en $x_1$, pero para que $h\cdot f\cdot\bar h$ sea un lazo basado en $x_1$ (con $f$ basado en $x_0$), $h$ debe comenzar en $x_1$ y terminar en $x_0$, que es lo que se usa después.}
\end{erratas}

\end{document}
